\subsection{Hessian spectra, profiling, and group heterogeneity}
The geometric conditioning and representation capacity of curvature must be evaluated together. The squared-distance calculation gives an exact conditioning factor for a tractable class of objectives. Profiling then accounts for the movement of the representation optimum, and the group-gap calculation quantifies a possible approximation benefit under explicit separability assumptions.
\begin{proposition}[Polar metric and distance Hessian]
\label{prop:squared-hessian}
For $r=d_K(x,y)$ and $h_y(x)=r^2/2$,
\begin{align}
 g_K&=dr^2+\frac{\sinh^2(\sqrt Kr)}K\,g_{\mathbb S^{d-1}},\\
 \operatorname{Hess}r
 &=\sqrt K\coth(\sqrt Kr)(g_K-dr\otimes dr),\\
 \operatorname{Hess}h_y
 &=dr\otimes dr+\sqrt Kr\coth(\sqrt Kr)(g_K-dr\otimes dr),\\
 \sigma(\operatorname{Hess}h_y)
 &=\{1,\sqrt Kr\coth(\sqrt Kr),\ldots,\sqrt Kr\coth(\sqrt Kr)\}.
\end{align}
If $d_K(x,y)\le D$,
\begin{align}
 I\preceq\operatorname{Hess}h_y\preceq\zeta_K(D)I,
 \qquad\zeta_K(D):=\sqrt KD\coth(\sqrt KD).
\end{align}
\end{proposition}
\begin{proof}
Let $a(r)=\sinh(\sqrt Kr)/\sqrt K$ and $E$ be a unit angular vector.
\begin{align}
 g_K(E,E)&=1,&\nabla_E\partial_r&=\frac{a'(r)}{a(r)}E,\\
 \operatorname{Hess}r(E,E)&=\frac{a'}a=\sqrt K\coth(\sqrt Kr),\\
 \operatorname{Hess}r(\partial_r,\cdot)&=0,\\
 \operatorname{Hess}(r^2/2)&=dr\otimes dr+r\operatorname{Hess}r,\\
 \frac d{dz}(z\coth z)
 &=\frac{\sinh z\cosh z-z}{\sinh^2z}\ge0,\\
 \sinh z\cosh z-z&=\int_0^z2\sinh^2s\,ds\ge0,\\
 z\coth z&=1+z^2/3-z^4/45+O(z^6).
\end{align}
\end{proof}

\begin{corollary}[Fr\'echet-type objectives]
For $f_K(x)=\mathbb E_y[d_K(x,y)^2/2]+R(x)$ with
$\mu_RI\preceq\operatorname{Hess}R\preceq L_RI$ and
$d_K(x,y)\le D$,
\begin{align}
 (1+\mu_R)I&\preceq\operatorname{Hess}f_K
 \preceq(\zeta_K(D)+L_R)I,\\
 \kappa(\operatorname{Hess}f_K)
 &\le\frac{\zeta_K(D)+L_R}{1+\mu_R},\\
 d_K(x,o_K),d_K(y,o_K)\le R&\Longrightarrow D\le2R.
\end{align}
\end{corollary}
\begin{proof}
\begin{align}
 \operatorname{Hess}\mathbb E_yh_y
 &=\mathbb E_y\operatorname{Hess}h_y,\\
 d_K(x,y)&\le d_K(x,o_K)+d_K(o_K,y).
\end{align}
\end{proof}

\begin{proposition}[Asymptotic growth matching]
For $d>1$, $b>1$, $\Delta r>0$,
\begin{align}
 \operatorname{Area}_K(r)
 &=\omega_{d-1}\left(\frac{\sinh(\sqrt Kr)}{\sqrt K}\right)^{d-1},\\
 \lim_{r\to\infty}\frac1r\log\operatorname{Area}_K(r)
 &=(d-1)\sqrt K,\\
 \lim_{\ell\to\infty}\frac1{\ell\Delta r}\log b^\ell
 &=\frac{\log b}{\Delta r},\\
 (d-1)\sqrt{K_{\rm match}}=\frac{\log b}{\Delta r}
 &\Longleftrightarrow
 K_{\rm match}=\left(\frac{\log b}{(d-1)\Delta r}\right)^2.\label{eq:capacity-heuristic}
\end{align}
\end{proposition}
\begin{proof}
\begin{align}
 \log\sinh z&=z-\log2+\log(1-e^{-2z}),\\
 \frac{\log\operatorname{Area}_K(r)}r
 &=\frac{\log\omega_{d-1}}r+
 (d-1)\left[\sqrt K-\frac{\log(2\sqrt K)}r+
 \frac{\log(1-e^{-2\sqrt Kr})}r\right].
\end{align}
\end{proof}

Growth matching identifies an asymptotic capacity scale; it does not by itself minimize a finite-data training loss. To study the loss, we instead follow a locally regular branch of representation optima. The Schur complement below shows that allowing representations to readjust subtracts a positive semidefinite term from the curvature Hessian. Consequently, positive curvature-block curvature at fixed representations need not imply a convex profiled problem.

\begin{assumption}[Profile regularity]
For every $u$ in an open set, let
\begin{align}
 q^*(u)&\in\arg\min_qF(q,u),&
 \operatorname{grad}_qF(q^*(u),u)&=0,\\
 H_{qq}(u)&:=\operatorname{Hess}_{qq}F(q^*(u),u)\succeq\mu I,&\mu&>0,\\
 \|H_{qu}(u)\|&\le B,&H_{uu}(u)&\succeq-L_{uu}I.
\end{align}
The branch $q^*$ is locally $C^1$ and $F$ is $C^2$.
\end{assumption}
\begin{proposition}[Implicit optimum derivative and Schur complement]
\label{prop:rigid-profile}
The envelope theorem \citep{milgrom2002envelope} and the implicit function theorem give, for $\phi(u)=F(q^*(u),u)$,
\begin{align}
 D_uq^*(u)&=-H_{qq}(u)^{-1}H_{qu}(u),\\
 \|D_uq^*(u)\|&\le B/\mu,\\
 \nabla\phi(u)&=\nabla_uF(q^*(u),u),\\
 \nabla^2\phi(u)&=H_{uu}-H_{uq}H_{qq}^{-1}H_{qu},\\
 \nabla^2\phi(u)&\succeq-(L_{uu}+B^2/\mu)I,\\
 \nabla^2[\phi(u)+\lambda\|u-u_0\|^2]
 &\succeq(2\lambda-L_{uu}-B^2/\mu)I.
\end{align}
\end{proposition}
\begin{proof}
At $q^*(u)$ choose normal coordinates in $q$ and differentiate
$\operatorname{grad}_qF(q^*(u),u)=0$:
\begin{align}
 H_{qq}D_uq^*+H_{qu}&=0,\\
 D\phi[h]&=F_q[D_uq^*h]+F_u[h]=F_u[h],\\
 D^2\phi[h,h]
 &=\langle h,H_{uu}h\rangle-
 \langle H_{qu}h,H_{qq}^{-1}H_{qu}h\rangle\\
 &\ge-L_{uu}\|h\|^2-\mu^{-1}B^2\|h\|^2.
\end{align}
\end{proof}

\begin{proposition}[Sharp heterogeneity bounds by weighted projection]
\label{prop:group-gap}
Let $I$ be an interval, $w_g>0$, $\sum_gw_g=1$, and
$0<\mu_g\le f_g''\le L_g$ on $I$, with $u_g^*\in I$ and $f_g'(u_g^*)=0$.
Define
\begin{align}
 a_g&:=w_g\mu_g,&b_g&:=w_gL_g,\\
 A&:=\sum_ga_g,&B&:=\sum_gb_g,\\
 u_a&:=A^{-1}\sum_ga_gu_g^*,&u_b&:=B^{-1}\sum_gb_gu_g^*,\\
 V_a&:=\sum_ga_g(u_g^*-u_a)^2,&V_b&:=\sum_gb_g(u_g^*-u_b)^2,\\
 \Delta_{\rm grp}&:=\min_{u\in I}\sum_gw_gf_g(u)-\sum_gw_gf_g(u_g^*).
\end{align}
Then the sharp bounds are
\begin{align}\label{eq:sharp-group-gap}
 \tfrac12V_a&\le\Delta_{\rm grp}\le\tfrac12V_b,\\
 V_a&=A^{-1}\sum_{g<h}a_ga_h(u_g^*-u_h^*)^2.
\end{align}
Writing $\bar u=\sum_gw_gu_g^*$ and $V_*=\sum_gw_g(u_g^*-\bar u)^2$ recovers
\begin{align}
 \tfrac12(\min_g\mu_g)V_*\le\tfrac12V_a
 \le\Delta_{\rm grp}\le\tfrac12V_b\le\tfrac12(\max_gL_g)V_*.
\end{align}
Both endpoints are attained by admissible quadratic profiles; if $\mu_g=L_g$ for every group, the lower and upper bounds coincide exactly.
\end{proposition}
\begin{proof}
Taylor's integral formula at $u_g^*$ gives
\begin{align}
 \tfrac12\mu_g(u-u_g^*)^2
 &\le f_g(u)-f_g(u_g^*)\le\tfrac12L_g(u-u_g^*)^2.
\end{align}
The weighted projection identity yields
\begin{align}
 \sum_ga_g(u-u_g^*)^2&=A(u-u_a)^2+V_a,\\
 AV_a&=A\sum_ga_g(u_g^*)^2-\left(\sum_ga_gu_g^*\right)^2
 =\sum_{g<h}a_ga_h(u_g^*-u_h^*)^2.
\end{align}
Minimizing the lower quadratic and evaluating the upper quadratic at $u_b\in I$ prove Eq.~\eqref{eq:sharp-group-gap}. Finally,
\begin{align}
 V_a&=\min_u\sum_gw_g\mu_g(u-u_g^*)^2\ge(\min_g\mu_g)V_*,\\
 V_b&\le\sum_gw_gL_g(\bar u-u_g^*)^2\le(\max_gL_g)V_*.
\end{align}
The profiles $f_g(u)=\mu_g(u-u_g^*)^2/2+c_g$ attain the lower endpoint; replacing $\mu_g$ by $L_g$ attains the upper endpoint.
\end{proof}

The bound is a weighted distance to the shared-curvature subspace. Keeping the individual Hessian bounds weights disagreement by each group's sensitivity, rather than charging every group the smallest $\mu_g$ or largest $L_g$. The pairwise formula shows which incompatible group preferences contribute to the approximation cost, and the quadratic equality cases establish exactly which constants are sharp.

\begin{corollary}[Exact quadratic profile]
For $f_g(u)=a_g(u-u_g^*)^2/2+c_g$, $a_g>0$,
\begin{align}
 u_{\rm shared}^*&=\frac{\sum_gw_ga_gu_g^*}{\sum_gw_ga_g},\\
 \Delta_{\rm grp}&=\frac12\sum_gw_ga_g(u_g^*-u_{\rm shared}^*)^2.
\end{align}
\end{corollary}
\begin{proof}
\begin{align}
 \frac d{du}\sum_gw_gf_g(u)&=u\sum_gw_ga_g-\sum_gw_ga_gu_g^*,\\
 \frac {d^2}{du^2}\sum_gw_gf_g(u)&=\sum_gw_ga_g>0.
\end{align}
\end{proof}
